Let $p(m)$ denote the least prime factor of $m$ (for $m\ge 2$), and interpret the quantifier ``$m$ is composite'' as part of the condition (otherwise the auxiliary quantity $F(n)=\min_{m>n}(m-p(m))$ would be identically $0$ because primes have $p(m)=m$).

Question: Does there exist a function $f:\mathbb{N}\to\mathbb{N}$ with $f(n)\to\infty$ such that for all sufficiently large integers $n$ there exists a \emph{composite} integer $m$ satisfying
\[
 n+f(n)<m<n+p(m)?
\]
Equivalently, writing $m=n+k$, we ask whether for all large $n$ there exists an integer $k>f(n)$ such that $n+k$ is composite and $p(n+k)>k$.

In terms of $F(n):=\min\{m-p(m): m>n\ \text{composite}\}$, the existence of such an $f$ is equivalent to the statement that $n-F(n)\to\infty$ (since $m<n+p(m)\iff m-p(m)<n$).
